% TOP_PROBLEM: {"id": 30006901, "title": "Margulis–Platonov normal-subgroup conjecture", "category_id": 17, "category": {"id": 17, "name": "group_theory", "display_name": "Group Theory", "description": "Problems about groups, group actions, representations, and related algebraic structures.", "slug": "group-theory", "order_index": 17, "created_at": "2026-07-31T15:26:25.671Z"}, "status": "open"}
% FIRST_POSTED: 2026-09-27
% !TeX program = lualatex
% ATTEMPT_STATUS: unresolved
\documentclass[11pt]{article}
\usepackage[margin=25mm]{geometry}
\usepackage{fontspec}
\setmainfont{Latin Modern Roman}
\usepackage{amsmath,amssymb,mathtools}
\usepackage{unicode-math}
\setmathfont{Latin Modern Math}
\usepackage{xurl,hyperref}
\hypersetup{colorlinks=true,urlcolor=blue}
\setcounter{secnumdepth}{0}
\setlength{\parindent}{0pt}
\setlength{\parskip}{5pt}
\setlength{\emergencystretch}{3em}
\begin{document}
\section*{\texorpdfstring{Margulis–Platonov normal-subgroup conjecture}{Margulis–Platonov normal-subgroup conjecture}}\label{30006901}
\subsection{Definitions and mathematical statement}
Let $\mathbf G$ be absolutely almost simple and simply connected over a number field $K$. For a place $v$, write $K_v$ for its completion and $\operatorname{rank}_{K_v}\mathbf G$ for the dimension of a maximal $K_v$-split torus. Let $A$ be the finite set of nonarchimedean places with rank zero and define
\[
 \delta:\mathbf G(K)\longrightarrow\prod_{v\in A}\mathbf G(K_v).
\]
The Margulis--Platonov conjecture asserts that every normal subgroup $N\triangleleft\mathbf G(K)$ not contained in the center has the form $N=\delta^{-1}(W)$ for an open normal subgroup $W$ of this product in its local-field topology. If $A$ is empty, the product is trivial, so the assertion says every noncentral normal subgroup is the whole group. It is not a classification only of normal subgroups of an arithmetic lattice.
\par\smallskip\textit{Formulation and concept references:} \href{https://arxiv.org/html/2603.27472}{[S1]}.

\subsection{Short English statement}
Are all noncentral normal subgroups of the rational points of a simply connected almost simple group obtained from open normal subgroups at its anisotropic finite places?
\subsection{Research attempt}
\paragraph{Scope and source audit (27 September 2026).}
The group being classified is $\mathbf G(K)$ itself. It is usually
neither discrete nor finitely generated in its local realizations.
A normal-subgroup theorem for an arithmetic lattice does not immediately
classify its abstract normal subgroups. The center is explicitly excluded
from the conclusion, and the local product contains only the specified
nonarchimedean anisotropic places.

Rapinchuk's 2026 primary paper [S1] states (MP) in precisely this form.
Its Theorem 1.1 \emph{assumes} (MP) to obtain a congruence-subgroup
conclusion for sets of places of positive density. The assumption must
not be read as a new proof of (MP). The introduction also records
established cases, including anisotropic inner forms of type $A$.
This attempt does not reprove those classification results or infer the
remaining cases from the density hypothesis of that different theorem.

The concrete work below gives a complete elementary proof of the target
for $\mathbf G=\operatorname{SL}_2$ over any number field. It then
formulates the topological factorization step needed for a nonempty
anisotropic product and tests why closing a subgroup is insufficient.
The split calculation illustrates a real mechanism, but does not assume
that anisotropic groups contain the unipotent subgroups it uses.

\paragraph{What the local description would immediately imply.}
Put $P=\prod_{v\in A}\mathbf G(K_v)$. Each local factor in this
product is compact: this is the standard compactness theorem for the
points of a semisimple anisotropic group over a nonarchimedean local
field. It is an external structural input, used here only to unpack
the proposed conclusion. The finite product $P$ is compact and totally
disconnected.

An open subgroup $W\le P$ has finite index. Its left cosets are open
and cover $P$, so compactness gives finitely many of them. If $W$ is
normal, the quotient $P/W$ is a finite group with discrete topology.
For $N=\delta^{-1}(W)$, the map
\[
 \mathbf G(K)/N\longrightarrow P/W,
 \qquad gN\longmapsto\delta(g)W
 \tag{394.1}
\]
is well-defined and injective. Therefore the asserted description
would force every noncentral normal subgroup to have finite index.
No density assertion about $\delta(\mathbf G(K))$ is needed for this
particular consequence.

The converse is stronger than finite index. It requires every relevant
finite quotient to be detected continuously by these particular local
factors. An arbitrary finite-index kernel is an algebraic object; its
openness for the pulled-back local topology is extra information.
When $A$ is empty there is no nontrivial local quotient at all, and
the assertion becomes the exact equality $N=\mathbf G(K)$.

\paragraph{Two-by-two elementary matrices over an infinite field.}
Let $F$ be an infinite field of characteristic zero; in particular it
may be our number field $K$. Define
\[
 u(t)=\begin{pmatrix}1&t\\0&1\end{pmatrix},\quad
 l(t)=\begin{pmatrix}1&0\\t&1\end{pmatrix},\quad
 h(a)=\begin{pmatrix}a&0\\0&a^{-1}\end{pmatrix}\quad(a\ne0).
\]
The subgroups $U=\{u(t):t\in F\}$ and $L=\{l(t):t\in F\}$ are
abelian. Direct multiplication gives
\[
 h(a)u(t)h(a)^{-1}=u(a^2t),\qquad
 [h(a),u(t)]=u((a^2-1)t).
 \tag{394.2}
\]
Here $[x,y]=xyx^{-1}y^{-1}$. Choosing any $a$ with $a^2\ne1$ shows
that every $u(t)$ is a commutator. The lower transvections are conjugate
to the upper ones, so they also lie in the derived subgroup.

For completeness, these transvections generate all of
$\operatorname{SL}_2(F)$. For $t\ne0$ set
\[
 w(t)=u(t)l(-t^{-1})u(t)
     =\begin{pmatrix}0&t\\-t^{-1}&0\end{pmatrix}.
\]
Then $w(t)w(-1)=h(t)$. If
$B=\left(\begin{smallmatrix}a&b\\c&d\end{smallmatrix}\right)$
has determinant one and $a\ne0$, Gaussian elimination gives
\[
 B=l(c/a)h(a)u(b/a).
 \tag{394.3}
\]
The lower right entry on the right is $(1+bc)/a=d$. If $a=0$, then
$c\ne0$, and multiplying on the left by $u(1)$ makes the upper left
entry nonzero; apply the preceding factorization and undo that
transvection. Thus $U,L$ generate the group. Equation (394.2) now
proves that $\operatorname{SL}_2(F)$ is perfect.

This proof uses the availability of field inverses and an element with
$a^2\ne1$. It is not a proof for matrices over an arbitrary ring, where
Gaussian elimination can fail, or for every small finite field.

\paragraph{A group-action simplicity lemma with its full proof.}
\textbf{Lemma 394.1.} Suppose a group $G$ acts faithfully and primitively
on a set $\Omega$. Fix $\omega\in\Omega$. Suppose its stabilizer
$G_\omega$ has an abelian normal subgroup $U$ whose conjugates generate
$G$, and suppose $G$ is perfect. Then $G$ is simple.

\emph{Proof.} Let $N\triangleleft G$ be nontrivial. The $N$-orbits
form a $G$-invariant partition: $g(Nx)=N(gx)$. Primitivity makes this
partition either the singleton partition or the single orbit $\Omega$.
The singleton case would make $N$ act trivially, contrary to faithfulness.
Thus $N$ is transitive.

For each $g\in G$ choose $n\in N$ with $n\omega=g\omega$. Then
$g=nh$ for some $h\in G_\omega$. In the quotient $G/N$, conjugation
by $g$ has the same action as conjugation by $h$. Since $U$ is normal
in $G_\omega$, the images of $gUg^{-1}$ and $U$ in $G/N$ coincide.
All conjugates of $U$ therefore have the same abelian image there.
They generate the quotient, so the quotient is abelian. Perfectness
of $G$ forces its every abelian quotient to be trivial, hence $N=G$.
\hfill$\square$

Primitivity is essential in this argument: transitivity of the original
action alone does not make the orbits of a normal subgroup either
singletons or all of $\Omega$. We will verify the stronger property
through double transitivity, rather than infer it without proof.

\paragraph{Apply the lemma to the projective line.}
Let $G=\operatorname{PSL}_2(F)$ act by fractional linear transformations
on $\mathbb P^1(F)=F\cup\{\infty\}$. The action of
$\operatorname{SL}_2(F)$ has kernel exactly its scalar center
$\{I,-I\}$: fixing $0$ and $\infty$ forces a diagonal matrix, and
fixing $1$ makes its diagonal ratio one. Thus the projective action
is faithful.

It is transitive because translations move points in $F$ and $w(1)$
interchanges $0$ and $\infty$. The stabilizer of $\infty$ contains
all translations $u(t)$ and hence acts transitively on $F$. Therefore
the action is doubly transitive. A doubly transitive action on a set
with at least three points is primitive: a block containing a fixed
point and one other point must, under its point stabilizer, contain
every other point; otherwise all blocks are singletons.

The upper triangular determinant-one matrices stabilize $\infty$,
and their conjugation preserves $U$, as follows either by multiplication
or by decomposing them into a diagonal matrix and a translation.
So the image of $U$ is abelian and normal in the stabilizer. Its
conjugates contain the image of $L$, and hence generate $G$ by
(394.3). Perfectness passes to quotients, so $G$ is perfect.
Lemma 394.1 proves that $\operatorname{PSL}_2(F)$ is simple.

\paragraph{Lift the conclusion back through the center.}
\textbf{Proposition 394.2.} Every noncentral normal subgroup of
$\operatorname{SL}_2(F)$ is the whole group.

\emph{Proof.} Let $N$ be such a subgroup and write $Z=\{I,-I\}$.
Its image in $\operatorname{PSL}_2(F)$ is a nontrivial normal subgroup,
so it is all of that simple group. Hence
$\operatorname{SL}_2(F)=NZ$. It follows that
$\operatorname{SL}_2(F)/N$ is generated by the image of the center
and is abelian. Since $\operatorname{SL}_2(F)$ is perfect, this quotient
is trivial.\hfill$\square$

The last perfectness step matters. Surjectivity onto a central quotient
alone only gives $NZ=G$; it does not imply $N=G$ for an arbitrary
central extension. For example in the direct product $H\times C_2$,
the subgroup $H\times1$ surjects onto the quotient by $1\times C_2$
without being the whole group. Our proof checks the extra condition
rather than discarding a possible central discrepancy.

For $\mathbf G=\operatorname{SL}_2$ over a number field, the diagonal
torus remains split of rank one at every completion. Thus $A$ is empty.
Proposition 394.2 is exactly the required (MP) conclusion in this
complete infinite family of instances. It also proves that the normal
closure of any noncentral rational matrix is the whole group, without
producing a finite bound on the number of conjugates needed for each
individual target matrix.

\paragraph{What survives in higher split matrix groups.}
For $n\ge3$, elementary matrices satisfy
\[
 [I+aE_{ij},I+bE_{jk}]=I+abE_{ik}
 \quad(i,j,k\text{ distinct}).
 \tag{394.4}
\]
The formula follows from $E_{ij}E_{jk}=E_{ik}$ and the vanishing of
the other relevant products. Gaussian elimination over the field again
shows that elementary matrices generate $\operatorname{SL}_n(F)$,
so (394.4) proves perfectness there as well.

If a normal subgroup already contains one nontrivial elementary
transvection, diagonal conjugation supplies arbitrary nonzero parameters
on that root subgroup, and conjugation by determinant-one signed
permutation matrices supplies the other positions. It then contains
all elementary matrices and is the whole group. This is a useful
conditional certificate, but it is not used here as a proof that an
arbitrary noncentral normal subgroup contains such a transvection.
That extraction requires its own argument. The complete result claimed
above remains the fully proved two-by-two case.

In a $K$-anisotropic group there may be no nontrivial $K$-defined
unipotent subgroups at all. The elementary-matrix mechanism therefore
cannot be copied merely by replacing matrix symbols with algebraic-group
notation. This is one concrete obstruction to generalizing the proof.

\paragraph{An exact topological extension lemma.}
Let $P$ be a profinite group and $D\le P$ a dense subgroup. Suppose
$\phi:D\to F$ is a homomorphism to a finite discrete group. Then
$\phi$ extends continuously to $P$ if and only if it is continuous
for the subspace topology on $D$.

\emph{Proof.} Necessity is immediate. For sufficiency continuity at the
identity supplies an identity neighborhood $U\subset P$ with
$D\cap U\subset\ker\phi$. In a profinite group choose an open normal
subgroup $V\subset U$. The map $D\to P/V$ is onto: every coset is
nonempty and open, so density makes it meet $D$. Since
$D\cap V\subset\ker\phi$, the map $\phi$ factors through a
homomorphism $P/V\to F$. Composing with $P\to P/V$ gives the desired
extension. Uniqueness follows because two continuous maps into the
Hausdorff finite set agreeing on a dense subset agree everywhere.
\hfill$\square$

Thus, in a setting where density of the rational image has separately
been established, continuity of the relevant finite quotient is enough
to construct an open normal $W$ giving its kernel. This lemma is stated
with density as an explicit hypothesis; no unproved approximation theorem
is slipped into it. It separates the topological extension step, which
is straightforward, from the arithmetic continuity statement that would
need proof for each remaining case of (MP).

\paragraph{Closing a normal subgroup does not recover the subgroup.}
A proposed general argument might take the closure of $\delta(N)$,
prove that it is open, and then assert
$N=\delta^{-1}(\overline{\delta(N)})$. Only the inclusion from left
to right is automatic. Even openness and normality of that closure
do not give the reverse inclusion.

A particularly elementary model has all the relevant finite-index
features. Fix distinct primes $p,q$, take $P=\mathbb Z_p$ additively,
$D=\mathbb Z$ and $N=q\mathbb Z$. Then $D$ is dense in $P$ and
$N$ has finite index $q$ in $D$. Since $q$ is a unit in $\mathbb Z_p$,
the set $q\mathbb Z$ is also dense in $P$. Thus
\[
 \overline N=P,\qquad D\cap\overline N=D\ne N.
\]
Equivalently the reduction map $\mathbb Z\to\mathbb Z/q\mathbb Z$
is not continuous for the $p$-adic topology: no subgroup $p^m\mathbb Z$
is contained in $q\mathbb Z$. All groups here are abelian, so normality
is no issue. This is not a counterexample to (MP), whose rational points
satisfy much stronger structural hypotheses. It pinpoints the false
abstract topological inference.

One also cannot replace the local product by arbitrary convenient
places. The entire point of (MP) is that the selected anisotropic
factors suffice. Adding factors may make continuity easier while
changing the assertion one has proved.

\paragraph{Lattice normal subgroups are a different problem.}
The subgroup $\operatorname{SL}_2(\mathbb Z)$ has proper noncentral
normal congruence subgroups, for example the kernel of reduction modulo
an integer $m>2$. Such a kernel is not a normal subgroup of
$\operatorname{SL}_2(\mathbb Q)$. Indeed
$u(m)$ lies in it, whereas conjugating by $h(1/m)$ gives $u(1/m)$
by (394.2), and this is not even an integral matrix.

This explicit calculation explains why a lattice congruence subgroup
does not contradict Proposition 394.2. It also demonstrates why a theorem
about normal subgroups of a finitely generated integral subgroup cannot
be transplanted to rational points without controlling conjugation by
all rational elements. Conversely, proving a statement about rational
points alone does not immediately classify all normal subgroups of
each arithmetic lattice.

\paragraph{Diagnostics and remaining gap.}
Exact rational matrix calculations check (394.2)--(394.4), the Gaussian
factorization on a finite grid, and the congruence-subgroup conjugation
diagnostic. The proof of simplicity uses the full projective-line action,
not a finite matrix enumeration. Residue calculations also check the
$\mathbb Z\subset\mathbb Z_p$ density mechanism modulo powers of $p$.

\textbf{Outcome: UNRESOLVED in general.} The selected conjecture is
proved here for split $\operatorname{SL}_2$ over every number field.
For a general group the missing step is a structural description of
noncentral normal subgroups strong enough to force the precise local
continuity and full preimage equality. Compactness, finite index,
closure, the elementary-matrix identities, and a congruence theorem
that assumes (MP) do not supply that step.

\subsection{Sources}
\begin{itemize}
\item[S1]A. S. Rapinchuk, The congruence subgroup property for S-arithmetic subgroups of simple algebraic groups when S has positive Dirichlet density, arXiv:2603.27472v1 (29 March 2026).
\url{https://arxiv.org/html/2603.27472}.
\textit{Formulation source.}
\end{itemize}
\end{document}
